\section{Introduction}\label{sec:introduction}

How many edges must one traverse to get from one vertex of a polytope to
another? In 1957 Hirsch proposed that a \(d\)-dimensional polyhedron with
\(n\) facets should require at most \(n-d\) steps. Santos disproved this
bound for bounded polytopes by constructing an example of dimension \(43\)
with \(86\) facets and two vertices at least \(44\) steps apart~\cite{santos}.
That example left open the possibility that some polynomial in \(n\)
bounds all such distances. We construct bounded polytopes whose diameter
grows faster than every fixed power of their facet count.

The diameter is the maximum, over all pairs of vertices, of the length of
a shortest path in the undirected vertex--edge graph. Our examples are
\emph{simple}, meaning that every vertex lies on exactly \(d\) facets.
We can even allow a move directly between any two vertices of a common
face of dimension at most \(s\). Write \(\operatorname{dist}_s(v,w)\)
for the minimum number of these moves, so that \(\operatorname{dist}_1\)
is the usual graph distance. For a positive function of \(n\), the
notations \(n^{o(1)}\) and \(n^{\omega(1)}\) mean that the ratio of
its logarithm to \(\log n\) tends to zero and to infinity, respectively.

\begin{theorem}[Superpolynomial diameter]\label{thm:main}
For each positive integer-valued \(s(n)=n^{o(1)}\) there is a family
of bounded simple rational polytopes \(P_n\), one for every sufficiently
large \(n\), with exactly \(n\) facets, dimension \(n^{o(1)}\), and
\(s(n)\)-move diameter \(n^{\omega(1)}\).

Moreover, there is a linear functional \(x\) on each \(P_n\), of positive
range \(\Delta x=\max_{P_n}x-\min_{P_n}x\), such that for every fixed
\(0<\alpha<1\),
\[
 \min_{\substack{v,w\ \mathrm{vertices\ of}\ P_n\\
               |x(v)-x(w)|\ge n^{-\alpha}\Delta x}}
       \operatorname{dist}_{s(n)}(v,w)=n^{\omega(1)}.
\]
\end{theorem}

Taking \(s(n)=1\) disproves the polynomial Hirsch conjecture. The objective
estimate explains where the long distances occur. On the same polytope,
every pair of vertices separated by even a fraction \(n^{-\alpha}\) of
the objective range requires a long walk. Intermediate vertices may take
any objective values.

There is also a bound with both the dimension \(d\) and the facet excess
\(E=n-d\) prescribed. This includes the classical \(d\)-step case
\(E=d\), as well as regimes in which the two parameters grow at very
different rates.

\begin{theorem}[Prescribed dimension and facet excess]\label{thm:excess}
For any \(a,b>0\) there exist \(c,\delta>0\) such that, for all
integers \(d,E\) with \(m=\min\{d,E\}\) sufficiently large, some
rational simple \(d\)-polytope with exactly \(d+E\) facets has two
vertices requiring at least
\[
                         cE\min\{E^a,d^b\}
\]
moves across faces of dimension at most \(\lfloor m^\delta\rfloor\).
\end{theorem}

At \(n=2d\), for every prescribed \(s(d)=d^{o(1)}\), there are
vertices at \(s(d)\)-move distance \(d^{\omega(1)}\) such that each
facet contains exactly one of them (Corollary~\ref{cor:spindle}).
In the polar these vertices correspond to two disjoint simplex facets
that together contain every vertex. A section between the facets yields
a Minkowski sum of two simplices with superpolynomial facet-adjacency
diameter (Corollary~\ref{cor:minkowski}).

\subsection{Diameter and linear programming}

The simplex method gives Hirsch's question an algorithmic interpretation.
Starting from a feasible vertex, a pivot can move to an adjacent vertex.
Even complete knowledge of the destination cannot shorten a walk below
the graph distance. General LP algorithms have other choices.
In the companion paper~\cite{lpcompanion}, a strongly polynomial algorithm
finds an optimum by recovering its primal and dual certificates from the
inequalities. Corollary~\ref{cor:packing} makes this distinction concrete. It gives
packing linear programs in which the origin and the all-ones vector are
the unique minimizer and maximizer of the sum of the coordinates, yet
every edge path between these explicit endpoints has superpolynomial length.

Dantzig recorded Hirsch's original question in his treatment of linear
programming~\cite[p.~160]{dantzig}. Klee and Walkup reduced the linear
conjecture to the case \(n=2d\) and gave an unbounded counterexample~\cite{kw}.
For bounded polytopes, Matschke, Santos and Weibel brought the dimension
of Santos's counterexample down to \(20\)~\cite{msw}. Davies and
coauthors constructed five-dimensional prismatoids whose width grows
linearly in their number of vertices~\cite{davies}. The polynomial
conjecture asks for a single polynomial valid in all dimensions
\cite[Conjecture~1.3]{santos}.

Larman's upper bound \(n\,2^{O(d)}\) shows that a superpolynomial
diameter forces \(d/\log n\) to tend to infinity~\cite{larman}.
Theorem~\ref{thm:main} places the examples within this range while
keeping \(d\) smaller than every positive power of \(n\).
The Kalai--Kleitman bound and its refinements give quasipolynomial upper
bounds as \(d\) varies~\cite{kalai,todd,sukegawa}. Liu, Ge and Ye
also obtain a quasipolynomial upper bound in the facet excess alone
\cite{liugeye}. Our construction concerns actual convex polytopes.
It complements work on combinatorial models of diameter, including
connected layer families~\cite{abstraction}. By polarity, the examples
also give rational simplicial polytopes whose facets have superpolynomial
diameter under adjacency across ridges~\cite{ziegler}.

\subsection{Why distances multiply}

The construction uses a linear coordinate \(x\) to measure progress through
a polyhedron. We call \(x\) \emph{time}, although a path can move in either
direction in \(x\). The elementary way to force a bend is to give
\(J+1\) independent coordinates \(z_i\) the lower bounds
\(z_i\ge |x-a|\). At a fixed \(x\), a point in a face of the slice of
dimension at most \(J\) must attain at least one bound. In particular,
some \(z_i=0\) when the path reaches \(x=a\). If an affine segment
crossed that level in its interior, its \(z_i\) coordinate would give a
line through the corner lying above both branches of \(|x-a|\).
Figure~\ref{fig:kink} shows why this is impossible. A convex piecewise
affine lower bound with many corners forces a separate bend at each one.

\begin{figure}[htbp]
\centering
\begin{tikzpicture}[x=1cm,y=.8cm,font=\small,>=Stealth]
  \fill[lagoon!9] (-2.9,2.63) -- (-2.9,2.45) -- (0,0)
       -- (2.9,2.45) -- (2.9,2.63) -- cycle;
  \draw[lagoon,thick] (-2.9,2.45) -- (0,0) -- (2.9,2.45);
  \draw[ink!65,dotted] (0,-1.05) -- (0,2.64);
  \draw[ember,thick,dashed] (0,0) -- (1.21,-1.02);
  \fill[ink] (0,0) circle[radius=2pt];
  \node[left=2pt,ink] at (0,-.2) {\(x=a\)};
  \node[ink] at (0,1.7) {feasible values of \(z_i\)};
  \node[lagoon,anchor=west] at (3.05,2.4) {\(z_i\ge |x-a|\)};
  \node[ember,anchor=west] at (1.32,-.91) {affine continuation};
\end{tikzpicture}
\caption{A coordinate that attains its lower bound at \(x=a\) cannot be affine
through the crossing. Continuing the left branch past \(a\) violates the
right-hand inequality. More generally, the two inequalities require an
affine line through the contact point to have slope at most \(-1\) on the
left and at least \(1\) on the right.}
\label{fig:kink}
\end{figure}


To make the cost grow faster than the number of inequalities, we place
a copy of an expensive crossing inside each of many disjoint intervals
of \(x\). Within one such interval, called an \emph{episode}, a pair
of auxiliary coordinates records progress by its difference \(w\).
A pair can be chosen for which \(w\) crosses from near zero to near one.
Several disjoint triples of smaller intervals are marked on the \(w\)-axis.
On a traversal of a triple, either part of the path describes an entire
inner crossing, with \(w\) as its time, or the path crosses an overlap
of two phases in \(x\). We put another copy of the test in each such
overlap. Figure~\ref{fig:episode} shows the two time coordinates.
The last overlaps in this recursion use the piecewise affine lower
bounds above.

The lower bound for the inner crossing is counted repeatedly. For a
sufficiently large crossing cost \(D\), one level with \(N\) episodes
and \(q\) triples per episode forces at least \(NqD/4\) segments. To count them
on an undirected walk, for each interval stop at the first hit of its
far endpoint and look back to the last hit of its near endpoint. The
selected subpaths occupy disjoint intervals of the \emph{path parameter}.
Even two visits to the same physical overlap can force different bends
(Figure~\ref{fig:return-traversals}).

The cost in facets is a separate matter. Giving every inner crossing
its own variables and inequalities would spend the very factor the
construction gained in distance. The copies instead use common clock
readings. Different copies on the same bank of clocks can contribute
tight inequalities only in disjoint reading intervals. Eliminating the
readings therefore requires combinations containing at most one row
from an inner copy. When a coordinate needed for that copy is missing,
a prescribed group of slacks can take its place. The scheduler is the
system of inequalities that makes the required slacks available at
the appropriate times. It must do this for many episodes using few
coordinates.

Section~\ref{sec:diameter} converts a finite crossing bound into
Theorems~\ref{thm:main} and~\ref{thm:excess}, including the passage
to bounded simple rational polytopes with an exact facet count.
Sections~\ref{sec:framework}--\ref{sec:recursion} establish the crossing
bound using the scheduler. Its numerical input is Theorem~\ref{thm:coarse},
which assigns states to successive times subject to a coordinatewise
inequality between neighbors. Sections~\ref{sec:numerical}--\ref{sec:finite}
construct these states and prove the inequality.
Sections~\ref{sec:layers}--\ref{sec:compiler} turn them into the scheduler's
linear constraints. Section~\ref{sec:geometry} gives the consequences
for complementary vertices, Minkowski sums and packing linear programs.
